\documentclass[11pt]{article}
\usepackage[margin=0.9in]{geometry}
\usepackage[colorlinks=true,urlcolor=blue]{hyperref}
\pagestyle{empty}
\begin{document}

\begin{flushleft}
Quang-Vinh Dang\\
British University Vietnam\\
Hung Yen, Vietnam\\
\href{mailto:vinh.dq4@buv.edu.vn}{vinh.dq4@buv.edu.vn}\\[12pt]
\today\\[12pt]
\end{flushleft}

\noindent To the Editors, \emph{Technological Forecasting and Social Change}

\vspace{10pt}
\noindent Dear Editors,

\vspace{6pt}
\noindent Please find enclosed our manuscript, \textbf{``From forecasting to real-time surveillance: An anytime-valid monitoring framework for technology diffusion, with 35 years of evidence from the United States,''} for consideration.

Technological forecasting has long served as a planning tool for anticipating how new technologies will diffuse, but the discipline's core method---fitting a growth curve to data after collection ends---offers no way to ask, in real time, whether a technology's diffusion has actually changed regime. That question has become practical rather than academic: high-frequency adoption series, from the U.S. Census Bureau's biweekly AI-use survey to firms' own internal telemetry, are now inspected continuously, and the classical statistics underlying diffusion forecasting invalidate under repeated inspection.

We propose reframing forecasting as continuous surveillance. The paper develops an anytime-valid monitoring framework combining (i) a causal measurement layer that treats survey and instrument redesigns as interventions on the observation process rather than genuine shifts in the underlying diffusion state, and (ii) one-sided mixture e-processes whose error control holds uniformly over time, so that continuous monitoring and data-dependent stopping do not invalidate the inference -- unlike the repeated significance tests that dashboards and forecasting practice currently rely on by default.

Applied to the Census Bureau's biweekly enterprise-AI-adoption series, the monitor declares a diffusion takeoff on 20 April 2025, correctly holds the widely discussed mid-2025 slowdown at strong but sub-threshold evidence that is subsequently vindicated by re-acceleration, and absorbs a 26-standard-error November 2025 question-wording redesign as a measurement event rather than a diffusion shock -- an event that conventional repeated testing misreads as the largest AI adoption shock on record. We further extend the surveillance logic across 35 years of U.S. digital diffusion, recovering the internet's 1995 takeoff, 2008 saturation onset, and e-commerce's 2004, 2018, 2020, and 2023 regime shifts as dated declarations, and place the intervening technology waves (data mining, cloud, big data, data science, deep learning) on the same timeline using canonical landmarks, since no official adoption series exist for them.

We believe the paper fits the journal's mission directly: it contributes a new method to technological forecasting practice and offers planners a live instrument rather than a retrospective curve. All data are public (U.S. Census Bureau BTOS and FRED series); analysis datasets and replication scripts reproducing every number, table, and figure are included as supplementary material. The manuscript is original, unpublished, and not under consideration elsewhere. We have no conflicts of interest to disclose.

Thank you for considering our manuscript. We look forward to your response.

\vspace{14pt}
\noindent Sincerely,\\[6pt]
Quang-Vinh Dang\\
British University Vietnam\\
Corresponding author

\end{document}
